%% Comandi e stili comuni a tutta la relazione %%

%% Colori: serie dei grafici (in ordine fisso) e inchiostri secondari %%
\definecolor{serie1}{HTML}{2A78D6}
\definecolor{serie2}{HTML}{EB6834}
\definecolor{serie3}{HTML}{1BAF7A}
\definecolor{serie4}{HTML}{EDA100}
\definecolor{inchiostro}{HTML}{52514E}
\definecolor{muto}{HTML}{898781}
\definecolor{griglia}{HTML}{E1E0D9}
\definecolor{asse}{HTML}{C3C2B7}
\definecolor{sfondocodice}{HTML}{F5F5F3}

%% Dati dei grafici, ricavati dai risultati dei benchmark da scripts/report_data.py %%
\newcommand{\dati}[1]{dati/#1.csv}

%% Nomi di file, funzioni e opzioni nel testo %%
\newcommand{\codice}[1]{\texttt{#1}}

%% Estratti di codice %%
\lstdefinelanguage{nasm}{
    morekeywords={mov,movsxd,shl,add,sub,neg,jnz,jl,jle,jmp,cmp,test,dec,ret,
        movaps,movss,movd,movdqa,xorps,subps,mulps,addps,haddps,pshufd,fld,
        vmovapd,vsubpd,vmulpd,vaddpd,vxorpd,vhaddpd,vextractf128,vzeroupper,
        vmovd,vpshufd,vinsertf128,vmovdqa},
    morekeywords=[2]{eax,ecx,edx,esp,esi,rdi,rsi,rdx,xmm0,xmm1,ymm0,ymm1,dword},
    morecomment=[l]{;},
    sensitive=false
}
\lstset{
    basicstyle=\ttfamily\footnotesize\linespread{1.0}\selectfont,
    keywordstyle=\color{serie1}\bfseries,
    keywordstyle=[2]\color{serie2},
    commentstyle=\color{muto}\itshape,
    stringstyle=\color{serie3},
    backgroundcolor=\color{sfondocodice},
    frame=single,
    rulecolor=\color{griglia},
    framesep=4pt,
    xleftmargin=4pt,
    xrightmargin=4pt,
    columns=fullflexible,
    keepspaces=true,
    showstringspaces=false,
    breaklines=true,
    tabsize=4,
    captionpos=b,
    literate={è}{{\`e}}1 {à}{{\`a}}1 {ò}{{\`o}}1 {ù}{{\`u}}1 {é}{{\'e}}1 {ì}{{\`i}}1
}
\renewcommand{\lstlistingname}{Listato}
\renewcommand{\lstlistlistingname}{Elenco dei listati}

%% Grafici: griglia e assi recessivi, linee sottili %%
\pgfplotsset{
    relazione/.style={
        width=0.95\textwidth,
        height=0.6\textwidth,
        axis line style={asse},
        tick style={asse},
        grid=major,
        major grid style={griglia, line width=0.4pt},
        tick label style={font=\footnotesize, color=inchiostro},
        label style={font=\small, color=inchiostro},
        legend style={font=\footnotesize, draw=none, fill=none, text=inchiostro},
        legend cell align=left,
        every axis plot/.append style={line width=1.2pt},
        mark size=2.2pt,
    },
}

%% Pseudocodice in italiano %%
\floatname{algorithm}{Algoritmo}
\renewcommand{\listalgorithmname}{Elenco degli algoritmi}
\algrenewcommand\algorithmicrequire{\textbf{Input:}}
\algrenewcommand\algorithmicensure{\textbf{Output:}}
\algrenewcommand\algorithmicif{\textbf{se}}
\algrenewcommand\algorithmicthen{\textbf{allora}}
\algrenewcommand\algorithmicelse{\textbf{altrimenti}}
\algrenewcommand\algorithmicfor{\textbf{per}}
\algrenewcommand\algorithmicforall{\textbf{per ogni}}
\algrenewcommand\algorithmicdo{\textbf{esegui}}
\algrenewcommand\algorithmicwhile{\textbf{finché}}
\algrenewcommand\algorithmicreturn{\textbf{restituisci}}
\algrenewcommand\algorithmicend{\textbf{fine}}
\algrenewcommand\algorithmicfunction{\textbf{funzione}}
\algrenewcommand\algorithmicprocedure{\textbf{procedura}}
\algrenewtext{EndIf}{\algorithmicend\ \algorithmicif}
\algrenewtext{EndFor}{\algorithmicend\ \algorithmicfor}
\algrenewtext{EndWhile}{\algorithmicend\ \algorithmicwhile}
\algrenewtext{EndFunction}{\algorithmicend\ \algorithmicfunction}
\algrenewtext{EndProcedure}{\algorithmicend\ \algorithmicprocedure}

%% Righe di tabella generate da scripts/report_data.py: \input non si può usare dentro tabular %%
\makeatletter
\newcommand{\tabella}[1]{\@@input tabelle/#1.tex }
\makeatother
